Road traffic violations are a major cause of accidents and fatalities worldwide. With the advancement of computer vision and deep learning technologies, automated traffic monitoring systems have gained significant attention for detecting violations such as helmetless riding, triple riding, mobile phone usage while driving, and vehicle identification through Automatic Number Plate Recognition (ANPR).

Recent studies have demonstrated the effectiveness of deep learning-based object detection models for traffic surveillance applications. Among various approaches, the YOLO (You Only Look Once) family of models has emerged as one of the most widely adopted techniques due to its high detection accuracy and real-time processing capability. Chen et al. proposed a deep learning framework for detecting helmet-rule violations among motorcyclists and reported promising results under complex traffic conditions \cite{Chen2024}. Similarly, Said et al. enhanced helmet detection performance using an improved PerspectiveNet architecture integrated with EfficientNetV2, achieving high accuracy in motorcycle rider classification tasks \cite{Said2024}.

Helmet violation detection has been extensively studied because the use of helmets significantly reduces the risk of fatal injuries in motorcycle accidents. Most existing approaches follow a two-stage pipeline in which motorcycles and riders are first detected, followed by classification of helmet usage. Recent advancements in object detection and image classification have enabled reliable helmet detection even in crowded traffic scenes \cite{Chen2024, Said2024}.

Another important traffic violation is triple riding, which poses a serious safety risk for motorcycle users. Deshpande et al. developed a deep-learning-based framework capable of identifying motorcycles, counting riders, and detecting triple-riding violations from surveillance footage \cite{Deshpande2025}. Their work demonstrated that rider-count estimation can be effectively integrated with object detection systems, although challenges such as occlusion, varying camera angles, and image quality continue to affect performance.

The detection of mobile phone usage while riding or driving has also become an active area of research. Distracted driving caused by mobile phone usage contributes significantly to road accidents. Charef et al. proposed a computer vision approach based on YOLOv8 for detecting mobile phone usage among motorcycle riders \cite{Charef2025}. Their findings highlighted the difficulty of accurately identifying handheld devices due to object size, rider movement, and partial occlusions. Nevertheless, recent advances in deep learning have shown promising results for real-time distraction detection systems.

Automatic Number Plate Recognition (ANPR) is a critical component of intelligent traffic management systems. ANPR systems automatically detect vehicle registration plates and extract textual information using image processing, optical character recognition (OCR), and deep learning techniques. Modern ANPR solutions commonly employ object detection models such as YOLO for plate localization and convolutional neural networks for character recognition, achieving high accuracy under varying environmental conditions \cite{Du2013}.

Several countries have integrated ANPR systems with electronic traffic enforcement mechanisms to automate violation management. E-Challan systems digitally generate traffic tickets based on detected violations and maintain centralized violation records. Such systems reduce manual intervention, improve transparency, and increase enforcement efficiency. Automated enforcement frameworks combining violation detection, ANPR, and digital ticket generation have been successfully deployed in various smart traffic management applications \cite{Sharma2023}.

Despite significant progress in individual components such as helmet detection, triple-riding detection, mobile phone detection, and ANPR, most existing studies focus on a single violation type. Comprehensive systems capable of detecting multiple violations and automatically generating digital challans remain relatively limited. Furthermore, the performance of these systems often depends on the availability of high-quality datasets that accurately represent real-world traffic conditions. Therefore, the development of an integrated multi-violation detection and automated e-challan system remains an important research area in intelligent transportation systems.